\small\begin{longtable}{@{}P{.09\linewidth}P{.59\linewidth}P{.24\linewidth}@{}}
\toprule Claim & Analytische uitspraak & Bron-ID\textquotesingle s \\\midrule\endfirsthead
\toprule Claim & Analytische uitspraak & Bron-ID\textquotesingle s \\\midrule\endhead
RC01 & MDM omvat technische en organisatorische beheertaken. & L05, L20 \\
RC02 & Dataproducten en federatieve benaderingen overlappen met MDM zonder ermee samen te vallen. & L12, L24 \\
RC03 & Reviewrapportagerichtlijnen ondersteunen transparantie, geen automatisch systematische dekking. & M01, M02 \\
RC04 & Masterdatadefinities benadrukken domeinentiteiten, hergebruik en afgestemd beheer. & B01, L18, L20 \\
RC05 & Een golden record is geen garantie van juistheid of institutioneel gezag. & L19, P01 \\
RC06 & Vroege MDM-publicaties bevatten semantisch en productgericht denken. & L12, L18, L19, L20, L24 \\
RC07 & Vier theoretische perspectieven bieden een eigen ordening van het corpus. & L07, L08, L15, L21 \\
RC08 & Datacategorieën verschillen in classificatiegrond en zijn niet allemaal exclusief. & L12, L20, S47 \\
RC09 & Capability-taxonomieën verschillen in granulariteit en doel. & L05, L20, P01 \\
RC10 & De lifecycle is een cyclisch analysekader, geen universeel stappenvoorschrift. & L20, P01, P02 \\
RC11 & Patroonlabels verschillen tussen MDM-bronnen. & L20, P01 \\
RC12 & Patroonlabels bepalen niet zelfstandig prestaties of consistentiegaranties. & L24, P01, P02 \\
RC13 & Entity resolution omvat meer dan tekstvergelijking. & L21, L29, L30 \\
RC14 & Identifiergelijkheid, recordgelijkheid en referentidentiteit verschillen. & L08, S38, S39 \\
RC15 & Matching, survivorship en domeinverandering vragen verschillende besluiten. & L08, L22, P01 \\
RC16 & ER-maatstaven kunnen verschillende rangordes opleveren. & L22 \\
RC17 & Datakwaliteit is contextgebonden en multidimensionaal. & L15, L19 \\
RC18 & Kwaliteitsdimensies vereisen expliciete vergelijkingsconventies. & L19, S12 \\
RC19 & Normscope, kwaliteitsmetadata en graafvalidatie hebben verschillende toetsobjecten. & S02, S03, S12, S32, S34 \\
RC20 & Governance adresseert besluitrechten en verantwoordelijkheid. & L07, L20, L27 \\
RC21 & Roltermen vragen bron- en contextgebonden interpretatie. & L07, L20, L23, S51 \\
RC22 & Standaardbeheer is geen volledige data governance. & L07, S23 \\
RC23 & Metadataregistratie onderscheidt conceptuele en representatieconstructies. & S07, S08, S09, S10 \\
RC24 & Een technisch schema bepaalt niet zelfstandig de volledige betekenis. & S08, S11, S14 \\
RC25 & Metadataregistries en catalogi hebben verschillende en overlappende scopes. & S08, S09, S47 \\
RC26 & Term, definitie en begrip zijn te onderscheiden; begrip/concept zijn bronafhankelijk. & S11, S13, S31 \\
RC27 & SKOS, NL-SBB en MIM hebben complementaire en overlappende functies. & S13, S14, S31 \\
RC28 & Bedrijfssemantiek komt aantoonbaar voor in vroegere MDM-literatuur. & B01, L20 \\
RC29 & Gelijke labels garanderen geen gelijke begripsafbakening. & B01, S31 \\
RC30 & EIF-niveaus verschillen van de eigen uitsplitsing van syntaxis en structuur. & S43 \\
RC31 & Conceptuele modellen veronderstellen ontologische categorieën en identiteit. & L08 \\
RC32 & Ontologische modelleerbenaderingen zijn geen universele MDM-vereiste. & L10, S40, S41 \\
RC33 & SKOS-concept, ontologische klasse, objecttype en data-element zijn geen automatische equivalenten. & S08, S13, S31 \\
RC34 & Semantic Web-specificaties vervullen verschillende representatie- en queryfuncties. & S28, S29, S30, S35, S36 \\
RC35 & SMDM is een bestaand antecedent van semantische MDM-integratie. & L18 \\
RC36 & Een formeel valide graaf garandeert geen juiste mapping of broninhoud. & L13, S32 \\
RC37 & PROV-O biedt constructies voor provenance; lineage is hier een werkafbakening. & S33 \\
RC38 & Herkomst is geen zelfstandige waarheids- of bevoegdheidswaarborg. & S25, S33 \\
RC39 & Valid time en transaction time zijn verschillende tijdsbetekenissen. & L14 \\
RC40 & Tijdnotatie en tijdontologie bepalen geen volledig historiebeleid. & P01, S37, S52 \\
RC41 & Documentsoort, normatieve status en empirisch bewijs zijn verschillende assen. & L11, S13, S30, S42 \\
RC42 & Cross-standard relaties vormen een synthese, geen verplichte componentenlijst. & S01, S02, S03, S07, S08, S09, S10, S11, S12, S13, S14, S15, S16, S17, S18, S19, S20, S21, S22, S23, S24, S25, S26, S27, S28, S29, S30, S31, S32, S33, S34, S35, S36, S37, S38, S39, S40, S41, S42, S43, S44, S45, S47, S49, S51, S52 \\
RC43 & Profielrelaties, complementaire metadata en onderzoekersmappings verschillen. & S16, S31, S33, S34, S45, S47 \\
RC44 & Overlap bewijst geen redundantie; scopebeperking bewijst geen wetenschappelijke lacune. & S14, S19, S31 \\
RC45 & Basisregistratieafspraken verbinden kerngegevens met institutionele verantwoordelijkheid. & S25, S26 \\
RC46 & BAG biedt begripsdocumentatie; geen operationele stelselketen is hier onderzocht. & S49 \\
RC47 & Nederlandse standaarden opereren op verschillende abstractieniveaus. & S13, S14, S15, S16, S17 \\
RC48 & Nederlandse lijststatus en technische versie moeten afzonderlijk worden vastgesteld. & F01, F02, F03 \\
RC49 & De FDS-dekking is te beperkt voor reconstructie van operationele afspraken. & S27 \\
RC50 & EIF, Europese wetgeving en SEMIC hebben verschillende functies. & S42, S43, S44 \\
RC51 & FAIR ondersteunt hergebruik en machinegericht gebruik zonder algemene openbaarheidseis. & L11 \\
RC52 & Masterdata en dataproducten beschrijven verschillende aspecten van gegevensaanbod. & L12 \\
RC53 & Masterdata als product is een expliciet eerder gepubliceerd idee. & L20 \\
RC54 & Productmetadata overlappen met catalogus- en registratiemetadata. & S08, S47, S51 \\
RC55 & Data Mesh-principes zijn beschreven in een grijze-literatuurreview. & L23 \\
RC56 & Federatie heeft technische en organisatorische dimensies. & L20, L23, L24 \\
RC57 & API- en eventstandaarden bepalen niet vanzelf domeinbetekenis. & S18, S19, S20, S21, S22 \\
RC58 & Data contracts en API-first zijn hier begrensde onderzoeksoriëntaties. & S19, S47, S51 \\
RC59 & AI kan hulpmiddel en consument van masterdata zijn; xLP is een beperkt voorbeeld. & L29, L31 \\
RC60 & RAG-resultaten bewijzen geen algemeen voordeel van semantische MDM. & L13, L16, S33 \\
RC61 & Machineleesbare epistemische status is een open effect- en representatievraag. & L16, S25, S33 \\
RC62 & Gegevens, representaties, beheer en consumptie vormen een eigen synthesemodel. & L07, L12, L19, L20, S08, S47 \\
RC63 & Gegevensrollen, correctheidsvormen en centralisatie-assen moeten worden onderscheiden. & L12, L22, L24, S32, S47 \\
RC64 & Convergentie in geselecteerde bronnen is geen vastgestelde vakgebiedbrede consensus. & L05, L07, L20, P01 \\
RC65 & Uniformiteit, autonomie, modelleerlast en bruikbaarheid vormen analytische spanningen. & L13, L19, L20, L23 \\
RC66 & De onderzoeksagenda onderscheidt synthesevragen van corpus- en toegangslacunes. & L05, L13, L14, L16, L18, L20, L21, L23, L24, L31, P01, S08, S13, S25, S33 \\
RC67 & Herbenoeming en technische vernieuwing blijven alternatieve interpretaties. & L13, L16, L18, L20 \\
RC68 & Integratieproblemen kunnen door ontbrekende toepassing in plaats van ontbrekende standaarden ontstaan. & L07, S08, S14 \\
\bottomrule\end{longtable}\normalsize
